\begin{solapartat}
\punts{0,3} Si apliquem balanç d'energia de l'efecte fotoelèctric tenim:
\[ E_C = hf - W_0 \]

\punts{0,2} La freqüència dels fotons és:
\[ f = \frac{c}{\lambda} = \frac{3,00\times 10^{8}}{300\times 10^{-9}} = 1,00\times 10^{15}\ \mathrm{Hz} \]

\punts{0,2} I l'energia dels fotons és:
\[ hf = 6,63\times 10^{-19}\ \mathrm{J}\,\frac{1\ \mathrm{eV}}{1,602\times 10^{-19}\ \mathrm{J}} = 4,14\ \mathrm{eV} \]

\punts{0,2} Per tant, la funció de treball que ens cal per construir el sensor és:
\[ W_0 = hf - E_C = 4,14 - 1 = 3,14\ \mathrm{eV}\,\frac{1,602\times 10^{-19}\ \mathrm{J}}{1\ \mathrm{eV}} = 5,03\times 10^{-19}\ \mathrm{J} \]

\punts{0,1} La freqüència llindar és:
\[ f_0 = \frac{W_0}{h} = 7,58\times 10^{14}\ \mathrm{Hz} \]

\punts{0,25} I finalment la longitud d'ona llindar és:
\[ \lambda_0 = \frac{c}{f_0} = \frac{3,00\times 10^{8}}{7,58\times 10^{14}} = 3,96\times 10^{-7}\ \mathrm{m}\,\frac{10^{9}\ \mathrm{nm}}{1\ \mathrm{m}} = 396\ \mathrm{nm} \]
\end{solapartat}

\begin{solapartat}
\punts{0,5} La freqüència llindar és:
\[ f_0 = \frac{W_0}{h} \]
I la longitud d'ona llindar s'expressa com:
\[ \lambda_0 = \frac{c}{f_0} = \frac{hc}{W_0} \]
Llavors pel tungstè
\[ \lambda_0 = h\,\frac{c}{W_0} = 6,63\times 10^{-34}\,\frac{3,00\times 10^{8}}{8,36\times 10^{-19}} = 2,38\times 10^{-7}\ \mathrm{m}\,\frac{10^{9}\ \mathrm{nm}}{1\ \mathrm{m}} = 238\ \mathrm{nm} \]
pel magnesi
\[ \lambda_0 = h\,\frac{c}{W_0} = 6,63\times 10^{-34}\,\frac{3,00\times 10^{8}}{5,86\times 10^{-19}} = 3,39\times 10^{-7}\ \mathrm{m}\,\frac{10^{9}\ \mathrm{nm}}{1\ \mathrm{m}} = 339\ \mathrm{nm} \]
i pel potassi
\[ \lambda_0 = h\,\frac{c}{W_0} = 6,63\times 10^{-34}\,\frac{3,00\times 10^{8}}{3,67\times 10^{-19}} = 5,42\times 10^{-7}\ \mathrm{m}\,\frac{10^{9}\ \mathrm{nm}}{1\ \mathrm{m}} = 542\ \mathrm{nm} \]

\begin{center}
\begin{tabular}{|l|l|l|l|}
\hline
Element & Símbol & Funció Treball (J) & Longitud d'ona llindar (nm) \\
\hline
Tungstè & W & $8,36\times 10^{-19}$ & 238 \\
\hline
Magnesi & Mg & $5,86\times 10^{-19}$ & 339 \\
\hline
Potassi & K & $3,67\times 10^{-19}$ & 542 \\
\hline
\end{tabular}
\end{center}

% DUBTE: la pauta escriu λ_{0,Mg} = 542 nm (hauria de ser λ_{0,K}) i "del tungstè i del potassi"
% (hauria de ser "del tungstè i del magnesi"); es transcriu tal com és.
\punts{0,75} Només podem fer servir el potassi que té una longitud d'ona llindar més gran que la que hem determinat a l'aparat a: $\lambda_{0,Mg} = 542\ \mathrm{nm} > 396\ \mathrm{nm}$. Les longituds d'ona llindars del tungstè i del potassi són massa petites.
\end{solapartat}
